% Compilar dos veces con pdflatex desde este directorio.
\documentclass[aspectratio=169,10pt]{beamer}
\usepackage[utf8]{inputenc}
\usepackage[T1]{fontenc}
\usepackage[spanish,es-nodecimaldot]{babel}
\usepackage{amsmath,amssymb,booktabs,tabularx,array,tikz}
\usetikzlibrary{arrows.meta,positioning,calc,shapes.geometric}

\definecolor{UAPNavy}{HTML}{17324D}
\definecolor{UAPTeal}{HTML}{007F82}
\definecolor{UAPGold}{HTML}{E2A33A}
\definecolor{UAPLight}{HTML}{F3F6F8}
\definecolor{UAPGray}{HTML}{536471}
\definecolor{UAPRed}{HTML}{A94343}

\tikzset{
  flow/.style={-{Stealth[length=2.2mm]},thick,draw=UAPNavy},
  datum/.style={draw=UAPGold,fill=UAPGold!16,rounded corners=2pt,align=center,minimum height=8mm,text=UAPNavy},
  process/.style={draw=UAPTeal,fill=UAPTeal!12,rounded corners=3pt,very thick,align=center,minimum height=10mm,text=UAPNavy},
  decision/.style={draw=UAPNavy,fill=UAPNavy!9,rounded corners=3pt,very thick,align=center,minimum height=10mm,text=UAPNavy},
  warning/.style={draw=UAPRed,fill=UAPRed!7,rounded corners=3pt,align=left}
}

\usetheme{Boadilla}
\usecolortheme{default}
\setbeamercolor{normal text}{fg=UAPNavy,bg=white}
\setbeamercolor{structure}{fg=UAPTeal}
\setbeamercolor{frametitle}{fg=white,bg=UAPNavy}
\setbeamercolor{title}{fg=white}
\setbeamercolor{subtitle}{fg=UAPGold}
\setbeamercolor{block title}{fg=white,bg=UAPTeal}
\setbeamercolor{block body}{fg=UAPNavy,bg=UAPLight}
\setbeamercolor{alertblock title}{fg=white,bg=UAPRed}
\setbeamercolor{alertblock body}{fg=UAPNavy,bg=UAPRed!7}
\setbeamerfont{frametitle}{series=\bfseries}
\setbeamertemplate{navigation symbols}{}
\setbeamertemplate{itemize item}{\color{UAPGold}$\blacktriangleright$}
\setbeamertemplate{itemize subitem}{\color{UAPTeal}$\bullet$}
\setbeamertemplate{footline}{%
  \leavevmode\hbox{%
    \begin{beamercolorbox}[wd=.78\paperwidth,ht=2.8ex,dp=1.2ex,leftskip=1em]{author in head/foot}%
      \color{white}\insertshorttitle
    \end{beamercolorbox}%
    \begin{beamercolorbox}[wd=.22\paperwidth,ht=2.8ex,dp=1.2ex,rightskip=1em plus1fil]{date in head/foot}%
      \color{UAPGray}\hfill Semana 7 \enspace | \enspace \insertframenumber/\inserttotalframenumber
    \end{beamercolorbox}}}

\newcommand{\concepto}[1]{\textcolor{UAPTeal}{\textbf{#1}}}
\newcommand{\pregunta}[1]{\begin{block}{Pregunta clave}#1\end{block}}
\newcommand{\br}{\\}
\title[IA · Clase 7]{Modelos supervisados fundamentales}
\subtitle{De funciones aprendidas a decisiones}
\author{Curso de Inteligencia Artificial}
\institute{Semana 7 · Clase 7}
\date{}

\begin{document}

% 1
\begin{frame}[plain]
  \begin{beamercolorbox}[wd=\paperwidth,ht=\paperheight,sep=2em]{frametitle}
    \vfill
    {\Huge\bfseries\inserttitle\par}\vspace{.45em}
    {\color{UAPGold}\Large\insertsubtitle\par}\vspace{1.2em}
    \begin{tikzpicture}[node distance=5mm,scale=.78,transform shape]
      \node[datum] (e) {Ejemplos};
      \node[process,right=of e] (f) {Función};
      \node[datum,right=of f] (p) {Predicción};
      \node[decision,right=of p] (a) {Decisión};
      \draw[flow] (e)--(f); \draw[flow] (f)--(p); \draw[flow] (p)--(a);
    \end{tikzpicture}\par
    \vfill {\color{UAPGold}\large\insertauthor\par}{\color{white}\insertinstitute\par}
  \end{beamercolorbox}
\end{frame}

% 2
\begin{frame}{Propósito de la clase}
  Comprender qué aprende un modelo supervisado y cómo su salida puede informar la acción de un agente.
  \begin{block}{Al finalizar, podremos}
    \begin{itemize}
      \item formular regresión y clasificación;
      \item explicar el ajuste lineal y la probabilidad logística;
      \item distinguir probabilidad, clase y acción;
      \item explicar cómo aprende un árbol;
      \item comparar modelos lineales, árboles y ensambles.
    \end{itemize}
  \end{block}
\end{frame}

% 3
\begin{frame}{Punto de partida}
  \small
  Ya estudiamos un agente que decide con una creencia conocida:
  \begin{center}\concepto{evidencia $\rightarrow$ probabilidad $\rightarrow$ costo esperado $\rightarrow$ acción}\end{center}
  Ahora la probabilidad o el valor esperado se estima desde ejemplos:
  \begin{center}\concepto{$(x_i,y_i)$ $\rightarrow$ modelo ajustado $\rightarrow$ predicción nueva}\end{center}
  \pregunta{¿Qué relación aprende cada familia de modelos?}
\end{frame}

% 4
\begin{frame}{Caso conductor: movilidad}
  \small
  Treinta minutos antes de una salida, un agente debe anticipar el estado de un trayecto.
  \begin{tabularx}{\textwidth}{>{\bfseries}p{.25\textwidth}X}\toprule
  Elemento & Ejemplo\\\midrule
  Entradas $x$ & duración planificada, congestión, lluvia y franja horaria\\
  Objetivo continuo & minutos de demora\\
  Objetivo binario & demora mayor que 10 minutos: sí/no\\
  Acción posterior & mantener, reasignar o solicitar revisión\\\bottomrule
  \end{tabularx}
  \begin{block}{Idea}El mismo dominio admite modelos distintos según la salida necesaria.\end{block}
\end{frame}

% 5
\begin{frame}{Aprendizaje supervisado}
  \small
  Disponemos de ejemplos con entrada y respuesta observada:
  \[\mathcal D=\{(x_1,y_1),\ldots,(x_n,y_n)\}.\]
  El algoritmo selecciona una función dentro de una familia:
  \[\widehat f=\arg\min_{f\in\mathcal F}\frac{1}{n}\sum_{i=1}^{n}L(y_i,f(x_i)).\]
  \begin{itemize}
    \item $\mathcal F$: funciones que el modelo puede representar.
    \item $L$: penalización por una predicción incorrecta.
    \item Aprender significa ajustar una función, no almacenar respuestas.
  \end{itemize}
\end{frame}

% 6
\begin{frame}{Tarea, modelo y decisión}
  \small
  \begin{tabularx}{\textwidth}{>{\bfseries}p{.20\textwidth}X}\toprule
  Nivel & Pregunta\\\midrule
  Tarea & ¿Qué salida necesitamos estimar?\\
  Modelo & ¿Qué familia de funciones utilizaremos?\\
  Ajuste & ¿Qué criterio determina sus parámetros?\\
  Decisión & ¿Cómo convertimos la salida en una acción?\\\bottomrule
  \end{tabularx}
  \begin{alertblock}{Distinción central}Un modelo predice; la regla de decisión incorpora costos, capacidad y restricciones.\end{alertblock}
\end{frame}

% 7
\begin{frame}{Regresión o clasificación}
  \small
  \begin{tabularx}{\textwidth}{>{\bfseries}p{.22\textwidth}p{.15\textwidth}p{.25\textwidth}X}\toprule
  Tarea & $y$ & Salida & Movilidad\\\midrule
  Regresión & real & $\widehat y$ & demora: 14,2 min\\
  Clasificación binaria & 0 o 1 & $\widehat p=P(Y=1\mid x)$ & probabilidad de demora crítica\\\bottomrule
  \end{tabularx}
  \begin{block}{Control}El puntaje probabilístico contiene más información que una clase rígida obtenida con un umbral.\end{block}
\end{frame}

% 8
\begin{frame}{Regresión lineal: hipótesis}
  \small
  La media del objetivo se aproxima mediante una función afín:
  \[\widehat y=\beta_0+\beta_1x_1+\cdots+\beta_px_p.\]
  \begin{itemize}
    \item $\beta_0$: valor de referencia.
    \item $\beta_j$: cambio asociado con una unidad de $x_j$, manteniendo las demás entradas fijas.
    \item “Lineal” se refiere a los parámetros; pueden incluirse transformaciones e interacciones.
  \end{itemize}
\end{frame}

% 9
\begin{frame}{Ejemplo de una recta}
  \small
  Horas de operación: $x=(2,4,6,8,10)$.

  Consumo en MWh: $y=(18,25,29,36,42)$.
  \[\widehat y=12{,}3+2{,}95x.\]
  Para siete horas:
  \[\widehat y=12{,}3+2{,}95(7)=32{,}95\text{ MWh}.\]
  \begin{block}{Lectura}La recta resume una tendencia dentro del rango observado.\end{block}
\end{frame}

% 10
\begin{frame}{¿Qué significa ajustar la recta?}
  \small
  El residuo del caso $i$ es $e_i=y_i-\widehat y_i$.

  Mínimos cuadrados elige los parámetros que minimizan:
  \[S(\beta)=\sum_{i=1}^{n}(y_i-\widehat y_i)^2.\]
  \begin{itemize}
    \item Los errores grandes pesan más porque se elevan al cuadrado.
    \item La pérdida define qué ajuste considera preferible el modelo.
    \item Ajustar no garantiza desempeño fuera de los ejemplos usados.
  \end{itemize}
\end{frame}

% 11
\begin{frame}{La solución en regresión simple}
  \small
  \[\widehat\beta_1=
  \frac{\sum_i(x_i-\bar x)(y_i-\bar y)}{\sum_i(x_i-\bar x)^2},
  \qquad \widehat\beta_0=\bar y-\widehat\beta_1\bar x.\]
  La pendiente es covarianza dividida por varianza de $x$.
  \begin{block}{Lectura del ejemplo}$2{,}95$ MWh adicionales por cada hora de operación dentro del rango estudiado.\end{block}
\end{frame}

% 12
\begin{frame}{Interpretar sin exagerar}
  \small
  Un coeficiente describe una asociación condicionada al modelo. No demuestra por sí solo que:
  \begin{itemize}
    \item modificar $x_j$ causará el cambio predicho;
    \item la relación continuará fuera del rango observado;
    \item las entradas omitidas son irrelevantes;
    \item el mismo coeficiente vale en otra población o período.
  \end{itemize}
  \begin{alertblock}{Límite}Asociación predictiva no equivale a efecto causal.\end{alertblock}
\end{frame}

% 13
\begin{frame}{Residuos: qué no explicó el modelo}
  \small
  Revisar residuos permite detectar límites de la forma elegida:
  \begin{itemize}
    \item curvatura: falta una relación no lineal;
    \item dispersión creciente: cambia la variabilidad;
    \item grupos separados: falta una diferencia relevante;
    \item observaciones extremas: pueden dominar el ajuste;
    \item patrón temporal: los errores no son independientes.
  \end{itemize}
  \begin{block}{Idea}Un error sistemático no se corrige agregando decimales.\end{block}
\end{frame}

% 14
\begin{frame}{Regresión lineal en movilidad}
  \footnotesize
  \[\widehat{\text{demora}}=-4+0{,}18\,\text{duración}+6{,}2\,\text{congestión}+3{,}1\,\text{lluvia}.\]
  Para duración 40 min, congestión 0,7 y lluvia 1:
  \[\widehat{\text{demora}}=-4+0{,}18(40)+6{,}2(0{,}7)+3{,}1=10{,}64.\]
  \begin{tabularx}{\textwidth}{>{\bfseries}p{.30\textwidth}X}\toprule
  Término & Lectura condicionada\\\midrule
  $0{,}18\,\text{duración}$ & 0,18 minutos por minuto planificado\\
  $6{,}2\,\text{congestión}$ & diferencia asociada con la escala de congestión\\
  $3{,}1\,\text{lluvia}$ & diferencia lluvia frente a no lluvia\\\bottomrule
  \end{tabularx}
\end{frame}

% 15
\begin{frame}{¿Por qué no usar una recta para probabilidades?}
  \small
  Si $Y\in\{0,1\}$, una recta puede producir valores menores que 0 o mayores que 1.
  \begin{block}{Necesitamos}
    \begin{enumerate}
      \item un puntaje que combine las entradas;
      \item una transformación que lo lleve al intervalo $[0,1]$.
    \end{enumerate}
  \end{block}
  La regresión logística utiliza un puntaje lineal y una función sigmoide.
\end{frame}

% 16
\begin{frame}{Regresión logística}
  \small
  Primero calcula un puntaje:
  \[\eta=\beta_0+\beta_1x_1+\cdots+\beta_px_p.\]
  Luego lo transforma:
  \[\widehat p=\sigma(\eta)=\frac{1}{1+e^{-\eta}}.\]
  \begin{itemize}
    \item $\eta$ puede tomar cualquier valor real.
    \item $\widehat p$ siempre pertenece a $(0,1)$.
    \item $\sigma(0)=0{,}5$ y $\sigma(-\eta)=1-\sigma(\eta)$.
  \end{itemize}
\end{frame}

% 17
\begin{frame}{Puntaje, probabilidad, clase y acción}
  \small
  \begin{tabularx}{\textwidth}{>{\bfseries}p{.22\textwidth}p{.15\textwidth}X}\toprule
  Objeto & Ejemplo & Función\\\midrule
  Puntaje $\eta$ & 0,8 & evidencia combinada\\
  Probabilidad $\widehat p$ & 0,69 & riesgo estimado\\
  Clase & crítica & salida de un umbral\\
  Acción & reasignar & decisión con costos y restricciones\\\bottomrule
  \end{tabularx}
  \begin{alertblock}{Control}Confundir estos objetos oculta dónde se produce el error.\end{alertblock}
\end{frame}

% 18
\begin{frame}{Ejemplo numérico de la sigmoide}
  \small
  Supongamos $\eta=-3+0{,}8x$.
  Para $x=4$:
  \[\eta=0{,}2,\qquad \widehat p=\frac{1}{1+e^{-0{,}2}}\approx0{,}55.\]
  Para $x=5$:
  \[\eta=1,\qquad \widehat p\approx0{,}731.\]
  \begin{block}{Resultado}El incremento de probabilidad no es constante: depende del punto de partida.\end{block}
\end{frame}

% 19
\begin{frame}{Odds y coeficientes}
  \small
  La regresión logística es lineal en los log-odds:
  \[\log\frac{p}{1-p}=\beta_0+\sum_j\beta_jx_j.\]
  Si $x_j$ aumenta una unidad, los odds se multiplican por $e^{\beta_j}$, manteniendo las demás entradas fijas.
  \begin{block}{Ejemplo}Si $\beta_j=0{,}40$, $e^{0{,}40}\approx1{,}49$: los odds aumentan 49\%, no la probabilidad en 49 puntos.\end{block}
\end{frame}

% 20
\begin{frame}{¿Cómo se ajusta la logística?}
  \small
  Para cada caso, $Y_i\sim\operatorname{Bernoulli}(p_i)$.

  El ajuste minimiza la log-loss:
  \[-\ell(\beta)=-\sum_i[y_i\log p_i+(1-y_i)\log(1-p_i)].\]
  \begin{itemize}
    \item Una probabilidad alta asignada a la clase equivocada recibe una penalización grande.
    \item El ajuste es iterativo: no existe la fórmula cerrada de mínimos cuadrados.
    \item La regularización puede limitar coeficientes inestables.
  \end{itemize}
\end{frame}

% 21
\begin{frame}{Probabilidad no significa certeza}
  \small
  Una salida $\widehat p=0{,}8$ significa 80\% de probabilidad bajo los datos y supuestos del modelo. No significa que:
  \begin{itemize}
    \item el evento ocurrirá;
    \item el modelo acierta 80\% en cualquier población;
    \item la probabilidad está bien calibrada;
    \item deba ejecutarse automáticamente una acción.
  \end{itemize}
  \begin{block}{Control}La calidad probabilística se comprueba con casos no usados para ajustar.\end{block}
\end{frame}

% 22
\begin{frame}{De probabilidad a acción}
  \small
  Decidir si actuar tiene dos errores posibles:
  \begin{itemize}
    \item actuar y que el evento no ocurra (falso positivo): costo $C_{FP}$;
    \item no actuar y que el evento ocurra (falso negativo): costo $C_{FN}$.
  \end{itemize}
  Se actúa cuando actuar es más barato que no actuar en valor esperado:
  \[(1-\widehat p)\,C_{FP}<\widehat p\,C_{FN}
  \qquad\Longleftrightarrow\qquad
  \widehat p>\frac{C_{FP}}{C_{FP}+C_{FN}}.\]
  Con $C_{FP}=8$ y $C_{FN}=40$: $t=\frac{8}{48}\approx0{,}167$.
  \begin{block}{Lectura}Si un falso negativo cuesta más que un falso positivo ($40>8$), conviene actuar incluso con probabilidades bajas: el umbral desciende por debajo de $0{,}5$.\end{block}
  \begin{alertblock}{Conclusión}El umbral 0,5 no es universal: depende de consecuencias, capacidad y restricciones.\end{alertblock}
\end{frame}

% 23
\begin{frame}{Caso conductor: probabilidad de demora}
  \small
  \[\eta=-3+0{,}04\,\text{duración}+2{,}2\,\text{congestión}+0{,}8\,\text{lluvia}.\]
  Para duración 40, congestión 0,7 y lluvia 1:
  \[\eta=-3+1{,}6+1{,}54+0{,}8=0{,}94,\]
  \[\widehat p=\sigma(0{,}94)\approx0{,}719.\]
  \begin{block}{Paso posterior}La regla de decisión determina si ese riesgo justifica intervenir.\end{block}
\end{frame}

% 24
\begin{frame}{Límite compartido de los modelos lineales}
  \small
  Regresión lineal y logística combinan entradas mediante una suma ponderada.
  \begin{columns}[T]
    \column{.48\textwidth}\begin{block}{Fortalezas}\begin{itemize}\item interpretación;\item cálculo eficiente;\item estabilidad con pocos datos.\end{itemize}\end{block}
    \column{.48\textwidth}\begin{alertblock}{Límite}\begin{itemize}\item frontera lineal;\item curvaturas explícitas;\item interacciones explícitas.\end{itemize}\end{alertblock}
  \end{columns}
\end{frame}

% 25
\begin{frame}[fragile]{Árbol de decisión}
  \small Un árbol aprende una secuencia de preguntas:
\begin{verbatim}
¿congestión > 0,65?
|-- sí: ¿llueve?
|   |-- sí  -> riesgo alto
|   `-- no  -> riesgo medio
`-- no      -> riesgo bajo
\end{verbatim}
  Cada hoja contiene una predicción basada en los ejemplos que llegan a ella.
\end{frame}

% 26
\begin{frame}{¿Cómo elige una división?}
  \small
  El árbol prueba atributos y umbrales y mide la impureza de cada nodo con el índice de Gini:
  \[G=1-p^2-(1-p)^2=2p(1-p),\]
  donde $p$ es la proporción de una de las clases en el nodo.
  \begin{itemize}
    \item $G=0$: nodo puro (una sola clase).
    \item $G=0{,}5$: clases equilibradas (máxima mezcla).
  \end{itemize}
  Se elige la división que más reduce la impureza, comparando la del padre con la suma ponderada de la de las hijas.
  \begin{block}{Lectura}$G$ es la probabilidad de que dos ejemplos tomados al azar del nodo tengan etiquetas distintas. Un padre $50/50$ ($G=0{,}5$) dividido en hijas $80/20$ y $20/80$ ($G=0{,}32$) deja la impureza ponderada en $0{,}32$, menor que la del padre.\end{block}
\end{frame}

% 27
\begin{frame}{Capacidad e inestabilidad del árbol}
  \small
  Un árbol profundo puede:
  \begin{itemize}
    \item representar relaciones no lineales e interacciones;
    \item aislar pocos ejemplos y memorizar ruido;
    \item cambiar ante pequeñas variaciones de los datos;
    \item producir probabilidades extremas en hojas pequeñas.
  \end{itemize}
  \begin{block}{Control}Profundidad, tamaño mínimo de hoja y poda regulan la complejidad.\end{block}
\end{frame}

% 28
\begin{frame}{Por qué combinar árboles}
  \footnotesize
  Los ensambles aprovechan que modelos distintos cometen errores diferentes.
  \begin{tabularx}{\textwidth}{>{\bfseries}p{.22\textwidth}p{.38\textwidth}X}\toprule
  Método & Mecanismo & Efecto\\\midrule
  Bagging & árboles sobre remuestras & reducir varianza\\
  Random Forest & remuestras y atributos aleatorios & aumentar diversidad\\
  Boosting & correcciones secuenciales & reducir error progresivamente\\\bottomrule
  \end{tabularx}
  \begin{alertblock}{Límite}Más modelos no reparan etiquetas incorrectas, información indebida o cambios de población.\end{alertblock}
\end{frame}

% 29
\begin{frame}{Random Forest}
  \small
  Un bosque aleatorio:
  \begin{enumerate}
    \item obtiene varias remuestras de los ejemplos;
    \item crece un árbol en cada remuestra;
    \item considera solo algunos atributos en cada división;
    \item promedia probabilidades o predicciones.
  \end{enumerate}
  \begin{block}{Intuición}El promedio estabiliza; la aleatoriedad evita que todos los árboles repitan las mismas primeras divisiones.\end{block}
\end{frame}

% 30
\begin{frame}{Boosting}
  \small
  Boosting construye un modelo aditivo:
  \[F_M(x)=F_0(x)+\sum_{m=1}^{M}\nu h_m(x).\]
  Cada árbol nuevo intenta corregir parte del error actual.
  \begin{itemize}
    \item $\nu$: tasa de aprendizaje.
    \item Árboles pequeños: correcciones acotadas.
    \item Muchas iteraciones: mayor capacidad y necesidad de control.
  \end{itemize}
  \begin{alertblock}{Diferencia}No es votación independiente: el aprendizaje es secuencial.\end{alertblock}
\end{frame}

% 31
\begin{frame}{Comparación de familias}
  \scriptsize
  \begin{tabularx}{\textwidth}{>{\bfseries}p{.19\textwidth}p{.25\textwidth}p{.22\textwidth}X}\toprule
  Aspecto & Lineal / logística & Árbol & Ensamble\\\midrule
  Relación & suma ponderada & reglas por regiones & combinación de regiones\\
  Interacciones & explícitas & automáticas & automáticas\\
  Interpretación & coeficientes & reglas si es pequeño & más difícil\\
  Extrapolación & posible, con riesgo & pobre & pobre\\
  Estabilidad & generalmente alta & baja & mayor que un árbol\\\bottomrule
  \end{tabularx}
  \begin{block}{Resultado}Ninguna familia domina en todos los problemas.\end{block}
\end{frame}

% 32
\begin{frame}{Cómo elegir un punto de partida}
  \small
  \begin{itemize}
    \item Salida numérica y relación aproximadamente suave: \concepto{regresión lineal}.
    \item Probabilidad binaria interpretable: \concepto{regresión logística}.
    \item Reglas no lineales e interacciones: \concepto{árbol controlado}.
    \item Datos tabulares con prioridad predictiva: \concepto{Random Forest o boosting}.
  \end{itemize}
  \begin{block}{Comparación justa}Mantener la misma tarea, información disponible y criterio de uso.\end{block}
\end{frame}

% 33
\begin{frame}{Lecturas}
  \small
  \begin{block}{Lectura principal}
    \begin{itemize}
      \item Capítulo 11, 11.1.1--11.1.5: regresión lineal.
      \item Capítulo 11, 11.4.1--11.4.5: regresión logística.
      \item Capítulo 12, 12.1.1--12.1.5: árboles.
    \end{itemize}
  \end{block}
  \begin{block}{Para ampliar}
    Capítulo 12, 12.3.1--12.3.5: ensambles, Random Forest y boosting.
  \end{block}
\end{frame}

\end{document}
